\chapter{Raw data and the log}

When a device answers strangely, the frames tell why. These two tabs show
what was really sent and received.

\section{The Raw data tab}
\label{sec:raw}

\screenshot[6cm]{raw-data}{A request and its answer on the Raw data tab.}
  {The Raw data tab with Decode fields on: the request (Data out) and its
   answer (Data in) field by field, and the history of frames below.}

The bytes exactly as sent to and received from the device, newest last:
\sent{Data out} in blue, \received{Data in} in green, as in the logo. Each
frame is decoded field by field (unit ID, function code, address, quantity,
byte count, data), and its check is verified: the CRC on RTU, the LRC on
ASCII, the header on TCP and UDP. A frame that fails its check is marked.

When no answer came, the tab says so (\emph{No response.}). \ui{Clear}
empties it.

\begin{tip}
To read a frame from another program or a bus analyser, paste it into the
Frame decoder (Section~\ref{sec:frame-decoder}): it decodes it the same way.
\end{tip}

\section{The Log tab}
\label{sec:log}

\screenshot[6cm]{log}{The Log tab, with three timeouts in red.}
  {The Log tab after a test with a few timeouts: time, function, address,
   quantity, duration and result of each request.}

Every request: \field{Time}, \field{Function},
\field{Address}, \field{Result} (OK, the exception, or a timeout) and
\field{Duration}. The tab's name counts them: \tab{Log (120)}.

\begin{description}[style=nextline, leftmargin=0pt, font=\intermedium]
  \item[\ui{Save log…}] Saves the log to a CSV file.
  \item[\ui{Clear log}] Empties it.
  \item[Right-click] \nav{Copy line}, \nav{Copy error}, \nav{Copy all}, and
    \nav{Show errors only} to keep just the failed requests in view.
  \item[Size] \nav{Settings\then Connection\then Requests kept in the log} sets how many
    requests are kept; older ones are dropped first.
\end{description}

For failures by cause over a long test, see Statistics
(Section~\ref{sec:statistics}).
